\documentclass[cameraready]{Interspeech}
\usepackage{microtype,array}
\title{Expert L2 English Speech Evaluation on a Gaming PC:\\Accuracy, Fluency, and Prosody with XLS-R, XEUS, and PhoneticXeus}
\author[affiliation={1}]{Richard}{Rose}
\address{$^1$HUFS, Republic of Korea}
\email{richard.rose@hufs.ac.kr}
\keywords{pronunciation assessment, fluency, prosody, speech representations, ridge regression}
\begin{document}
\maketitle
\begin{abstract}
\textit{oollim} is a locally executable system for evaluating L2 English pronunciation accuracy, fluency, and prosody on a consumer computer. The system combines pretrained XLS-R, XEUS, and PhoneticXeus representations with reference-text features and ridge regressions trained on SpeechOcean762, the Korean Learners' L2 Speech Corpus (English), J-AESOP, and their combinations. On 305 independently held-out recordings assessed by three expert raters, several ridge configurations substantially outperform archival Microsoft Azure Pronunciation Assessment API data in both Spearman correlation and quadratic weighted kappa across all three constructs. These results support expert speech evaluation on consumer hardware using pretrained representations and construct-specific regressions, without neural-model fine-tuning or calibration on the private evaluation data.
\end{abstract}

\section{Introduction}
Pronunciation evaluation requires distinct judgments of accuracy, fluency, and prosody. Accuracy concerns the realization of expected consonants and vowels \cite{deterding,cai}; fluency concerns speech continuity, including speed, pauses, and repairs \cite{kormos,bosker}; prosody concerns rhythm, stress, intonation, and phrasing \cite{cole,cutler}. These related dimensions inform pronunciation judgments but are not interchangeable \cite{cai,isaacs}. Segmental accuracy concerns distinctions between speech sounds, whereas intelligibility concerns what listeners understand and comprehensibility their perceived processing difficulty; a noticeable accent need not prevent understanding \cite{munro}. An accuracy score does not directly measure intelligibility. Fluency research distinguishes underlying processing efficiency, observable timing, and perceived smoothness \cite{segalowitz}. Speech rate, pause placement and duration, and repairs contribute differently to listener judgments, so fluency is not simply speaking quickly \cite{kormos,bosker}. Prosody organizes speech through lexical stress, prominence, phrasing, and intonation, with pitch, duration, and intensity providing interacting cues \cite{cole,cutler}. Appropriate stress can improve listeners' recall and evaluations of a speaker \cite{hahn}; contextual appropriateness therefore matters alongside acoustic variation. Although their acoustic correlates overlap, the three constructs require separate targets and useful measurements for educators; the scores need not come from the same representation combination or scoring function. XLS-R, XEUS, and PhoneticXeus provide complementary sources of speech information. This study examines whether these existing models, combined with modest supervised scoring functions, can support expert L2 English evaluation on an ordinary gaming PC. With computationally expensive pretraining supplied by the original model developers, the contribution is the selection and combination of their representations, reference-text evidence, and corpus-specific assessment supervision. The central evidence is an external evaluation on independently collected, privately held recordings, while training-corpus development results serve model selection and reproducibility and are provided in the repository. The external comparison includes Azure pronunciation assessment \cite{azure}, seven ridge-training families, and two human agreement baselines, with both rank association and numerical agreement measured. This design separates three questions: whether scores rank recordings similarly to human judgments, whether they agree on the numerical scale, and how that agreement compares with agreement among the human raters \cite{williamson}.

\section{Speech models and related work}
Automated assessment has long combined acoustic evidence with supervised mappings to human judgments. SpeechRater derives pronunciation, fluency, and other proficiency features from recognized speech and uses learned scoring models \cite{speechrater}. Neural assessment extends this approach by learning representations and their relationship to rated performance. GOPT uses phonetic posterior features and canonical phone identities in a transformer trained jointly on phone-, word-, and utterance-level human annotations \cite{gopt}. On SpeechOcean762, the 2,500-utterance training and test partitions provide the same assessment targets on different recordings \cite{speechocean,gopt}. The scoring model is trained on human ratings and evaluated against ratings of held-out recordings. Kim et al. adapt wav2vec 2.0 and HuBERT with CTC, then combine contextual representations and reference text in a bidirectional LSTM scoring model, evaluating against expert scores on separate partitions of Korean learner speech and SpeechOcean762 \cite{kim}. Cai et al. connect hierarchical neural assessment to applied-linguistics constructs, adapting scoring toward human judgments of segmental, suprasegmental, and temporal performance \cite{cai}. Saito et al. use regression over phonological, prosodic, and temporal features to predict comprehensibility, testing generalization on additional recordings and an interview task \cite{saito}. These studies train scoring models against human ratings and assess multiple dimensions or transfer to new data, although differences in targets, rating panels, and metrics prevent direct numerical comparison with \textit{oollim}. XLS-R extends cross-lingual self-supervised speech representation learning \cite{xlsr}; XEUS learns multilingual representations across thousands of languages \cite{xeus}; PhoneticXeus adds phone recognition and supplies hidden representations and phone evidence \cite{phonetic}. In \textit{oollim} their pretrained weights remain unchanged, while ridge regression learns regularized linear mappings to assessment labels \cite{ridge}. \textit{oollim} combines these representations with reference-phone diagnostics and construct-specific regressions trained on one or more corpora, then evaluates them on an independently collected passage-reading cohort to test transfer beyond holding out recordings within a training corpus. Earlier neural-head, temporal, and acoustic-feature experiments motivated the ridge baseline; the evaluation does not cover every neural assessment architecture. The comparison therefore focuses on the external transfer of compact scoring models.

\section{Method}
\subsection{Audio, reference text, and representations}
Whole-recording embeddings and phone diagnostics derived from required reference text provide the scoring inputs. FFmpeg \cite{ffmpeg} decodes audio-bearing containers; a common preprocessing path averages channels to mono and resamples to 16 kHz. The full recording is processed as one utterance, without averaging separately scored sentences. Each encoder adapter returns frame embeddings and valid-frame lengths or masks, with XLS-R 300M, XEUS, and PhoneticXeus each providing 1,024-dimensional frame representations in the configurations used here \cite{xlsr,xeus,phonetic}. Statistics pooling \cite{pooling}, using masked temporal means and population standard deviations, with $10^{-5}$ added to the variance before the square root, yields 2,048-dimensional utterance vectors. Model revisions and implementation checksums are pinned in the repository. A pinned CMU pronunciation dictionary supplies ARPAbet reference pronunciations \cite{cmudict}. ARPAbet--IPA correspondences are documented in CMU dictionary tools \cite{cmutools}; the versioned mapping in \textit{oollim} converts dictionary phones to the PhoneticXeus token inventory, including explicit diphthong splitting and unstressed-vowel handling. PhoneticXeus provides frame-level phone probabilities and a greedy connectionist temporal classification (CTC) sequence \cite{phonetic,ctc}. Seventeen diagnostics summarize phone probabilities and entropy \cite{shannon}, blank behavior, predicted and expected phone rates, reference coverage, relative sequence length, edit operations \cite{editdistance}, and CTC negative log-likelihoods normalized by frame and reference-token counts \cite{ctc}. The CTC calculation sums valid paths for the expected phone sequence rather than assigning a single forced alignment \cite{ctc}. These features combine pronunciation, timing, and uncertainty evidence; they are not error-free phonetic diagnoses. Out-of-vocabulary references are rejected explicitly. The evaluated accuracy configuration uses XEUS + PhoneticXeus + IPA diagnostics, whereas fluency and prosody use XLS-R + XEUS. The scoring pipeline can select a representation combination and training family separately for each construct.

\subsection{Training corpora and label mappings}
Three research corpora supply assessment labels mapped to a common 0--10 range. SpeechOcean762 supplies utterance accuracy, fluency, and prosodic ratings on the 0--10 scale. The official SpeechOcean762 training partition is divided by speaker into 2,000 training and 500 development utterances; the official test partition is excluded. The Korean Learners' L2 Speech Corpus (English) \cite{han} is distributed as AIHub dataset 71463 \cite{aihub}. Eligible Training PRN recordings yield 2,699 training and 636 development examples. The corpus rubric rates pronunciation accuracy (segment clarity and errors) and prosodic fluency (stress, rhythm, intonation, rate, and pauses) on 1--5 scales \cite{han}. The supplied articulation score trains accuracy, while the prosodic-fluency score trains both fluency and prosody after mapping with $2.5(y-1)$. J-AESOP \cite{aesop} supplies 180 Japanese L2 speakers reading the North Wind and Sun passage (task 6\_01). Deterministic speaker partitions allocate 108 recordings to training, 36 to development, and 36 to a reserved partition; native controls are excluded. Ratings from 16 raters across three sections are averaged per construct. Segment ratings supply accuracy, alongside fluency and prosody; $10(y-1)/9$ maps the 1--10 source scale to 0--10. Partitions are speaker-disjoint within each corpus to separate development from fitting \cite{cawley}. In the Korean Learners' L2 Speech Corpus (English), recorder identities linked by identical decoded audio are grouped before splitting; duration and dictionary-coverage checks determine eligibility. Exact sampling rules, audio audits, mappings, and development results are retained in the reproduction package. Equal numerical ranges do not establish score equivalence \cite{dorans}; the shared prosody/fluency target does not provide an independent fluency annotation. The mappings permit joint fitting while retaining these differences in supervision.

\subsection{Ridge fitting and family selection}
Seven training families compare single-corpus and mixed-corpus supervision. The single-corpus families are \texttt{speechocean}, \texttt{korean}, and \texttt{aesop}; \texttt{korean} abbreviates the Korean Learners' L2 Speech Corpus (English) in model identifiers. The pairs are \texttt{koreanocean}, \texttt{aesopocean}, and \texttt{koreanaesop}, while the three-corpus family is \texttt{aesopkoreanocean}. In the matched family procedure, a recording from corpus $c$ receives weight $w_i=1/(K n_c)$, where $K$ is the number of included corpora and $n_c$ their respective training counts, giving each corpus equal total weight regardless of size. Training-only weighted means and standard deviations standardize features \cite{esl}; each feature block (one encoder vector or the IPA diagnostics) of width $d_b$ receives gain $\sqrt{2048/d_b}$ to balance contributions from blocks of different sizes. Coordinates with standard deviation below $10^{-8}$ use scale 1. Each construct fits a weighted ridge objective \cite{ridge,esl}
\begin{equation}
\min_{\beta,b}\quad\sum_i w_i(y_i-z_i^\top\beta-b)^2+\alpha\lVert\beta\rVert_2^2,
\end{equation}
where $z_i$ is the standardized, block-scaled feature vector, $y_i$ the target score, $\beta$ the scoring coefficients, $b$ the unpenalized intercept, and $\alpha$ the regularization strength. Penalties $\{0.001,0.01,0.1,1,10,100\}$ are compared using an equally weighted mean of development Spearman correlations over the included corpora, separately for each construct; ties prefer lower mean MAE and then larger regularization, keeping selection on development data \cite{cawley}. The three regressions have separate coefficients, even when stored in one multi-output array. All seven training families are compared under three regularization-selection rules, holding corpus partitions, feature scaling, and the fitting procedure fixed across rules. The first rule maximizes development Spearman correlation separately for each construct. The two alternative rules minimize development accuracy mean absolute error (MAE), separately for each family and feature combination. Fluency and prosody share the accuracy-selected penalty for their common XLS-R + XEUS inputs, while their scoring coefficients remain separate; accuracy uses its own feature combination. One condition searches the same candidate grid as the Spearman condition; the other uses the historical grid, rescaled to preserve equivalent regularization under normalized training loss (see supplement). SpeechOcean is refitted under each rule for a matched comparison. The main table reports construct-specific leaders across all three rules; the supplement and aggregate release retain the full comparison and the original deployed SpeechOcean results. Public-development selection retained koreanocean for accuracy and aesopkoreanocean for fluency and prosody; these ranking-oriented research selections are distinct from the column maxima in the external table and are not claimed to maximize external QWK. The representation sensitivity analysis is separate from the fixed-combination external evaluation. A sensitivity check of all 315 family--construct--representation configurations on public development data retained the overall construct-specific selections under the 0.01 Spearman improvement threshold and corpus-level degradation and uncertainty checks specified before fitting in the representation-ablation protocol; the full protocol and results are provided in the reproduction materials. All coefficient estimation, preprocessing fitting, and score calibration exclude private evaluation ratings.



\subsection{Local execution and capacity}
Local operation was verified on consumer hardware using CPU inference. On the Windows gaming PC (Intel Core i7-11700K, 32 GB physical RAM, Ubuntu under WSL2 with approximately 24 GiB available RAM and 8 GiB swap), an archived 120-second capacity test completed in approximately 169 seconds, reaching 22.52 GiB peak process resident memory and 5.51 GiB peak system-wide swap use. The installed RTX 3080 was not required for that inference test. A separate Apple Silicon installation test processed an 8.21-second public recording in 42.13 seconds with four CPU threads, including decoding, all three encoders, phone features, and scoring. Software versions for the same WSL installation are recorded in the supplement. These single-run checks establish operation rather than representative throughput; the released duration limit is 120 seconds, and long recordings can require substantial RAM and swap.

\section{External evaluation}
\subsection{Independent recordings and ratings}
The collection comprised university EFL learners from China ($n=98$), Korea ($n=101$), and Japan ($n=107$), who recorded the same ``Please Call Stella'' passage on personal smartphones. One recording was excluded because of its length, leaving 305 samples. The three university professors had approximately 20--30 years of teaching experience, including pronunciation instruction, Cambridge English speaking assessment, and high-stakes speech rating. Based in those countries, the raters calibrated their judgments on separate recordings spanning proficiency levels, discussing discrepancies to establish shared criteria for accuracy, fluency, and prosody. Each professor then rated every recording independently in randomized order, blinded to participant identity and automated scores, using a 0--100 scale in five-point increments. The private evaluation recordings and ratings were held out from model training, hyperparameter tuning, and score calibration. The original human ratings and Azure API scores are archival; the seven ridge families were evaluated for the present study against those existing ratings. Azure scores were obtained in January 2026 using the \texttt{en-US} locale, with the Stella reference text supplied and each complete recording assessed as a whole. The returned accuracy, fluency, and prosody scores were retained. The service revision and remaining request options were not retained in the supplied evaluation archive; the comparison therefore concerns those recorded API outputs rather than a newly standardized Azure benchmark.

\subsection{Metrics and human baselines}
Evaluation distinguishes rank consistency from numerical agreement \cite{stemler,williamson}. Spearman's $\rho$ measures rank association \cite{spearman}. Quadratic weighted kappa (QWK) measures chance-corrected agreement with squared category-distance penalties \cite{cohen}; this study uses fixed integer categories 0--100 and rounding to the nearest integer, with exact half-points rounded upward. Model outputs are multiplied by ten; no calibration is fitted. Predictions are clipped to 0--100 for QWK, including bootstrap replicates, whereas Spearman retains raw predictions. A reference mean requires all three human scores; partial means are not substituted. Human--Human pairwise agreement averages the three pairwise statistics, whereas Human--Human consensus agreement averages each rater's agreement with the mean of the other two, never including that rater in their own reference. For each system, the three-rater row compares the system's scores with the mean of all humans, while the two-rater row averages the three separate agreement statistics against the three possible two-human means. Statistics are averaged, not observations or reference means pooled. The two-rater system row and Human--Human consensus row therefore provide the appropriate reference-size comparison when recomputed on the same recordings. Percentile 95\% intervals resample speaker clusters jointly, recomputing all component statistics within each replicate \cite{efron,field}. The supplementary marginal intervals retain the supplied evaluation export (2,000 replicates, seed 42). Additional paired comparisons use 600,000 bootstrap replicates per seed (42 and 2026), subtracting statistics on the same resampled recordings. Each unique recording URL is treated as one speaker. Bonferroni-adjusted intervals cover all 258 differences: 21 S/D/H models versus human consensus and Azure, plus Azure versus human consensus, for three constructs and two metrics. Both seeds agree on every adjusted interval-exclusion decision. These post-evaluation comparisons condition on the observed raters and fitted models; adjustment covers the reported contrasts, not the preceding model search. The supplement reports two-rater comparisons and intervals for the same 305 recordings; the aggregate release also includes three-rater and human-pairwise results. Matched reference sizes and cohorts make the human and model estimates comparable.



\section{Results}
Table~\ref{tab:external} reports observed leaders and near-ties under the .001 threshold. Rank association and numerical agreement favor different scorers: aesopocean leads accuracy correlation, koreanaesop and aesopkoreanocean supply fluency correlation leaders, and koreanaesop supplies all prosody correlation leaders. Koreanocean leads QWK, using S for accuracy and H for fluency and prosody. Across S/D/H, paired comparisons support 61 of 63 correlation advantages and 57 of 63 QWK advantages over Azure after Bonferroni correction across 258 contrasts. For all three constructs, the displayed Spearman leaders significantly exceed both the individual-human and archival Azure baselines on matched two-rater references, after Bonferroni correction. Against the individual-human baseline on matched two-rater references, S-koreanocean achieves significantly higher accuracy QWK than the human baseline (0.821 versus 0.704; Bonferroni-adjusted difference interval [0.051, 0.187]). H-koreanocean achieves slightly higher fluency QWK (0.744 versus 0.739), but the difference is not statistically significant. H-koreanocean's prosody QWK remains significantly below the human baseline. Both bootstrap seeds agree on all adjusted decisions. The observed leaders remain descriptive external selections.

\begin{table*}[t]
\centering
\small
\setlength{\tabcolsep}{5pt}
\renewcommand{\arraystretch}{1.0}
\caption{Construct-specific leaders and ties ($N=305$).}
\label{tab:external}
\begin{tabular}{@{}lllcccc@{}}
\toprule
Construct & Family & Rule & $\rho$ & 95\% CI & QWK & 95\% CI\\
\midrule
Accuracy & \textbf{Human consensus} & -- & \textbf{0.809} & [0.769, 0.842] & \textbf{0.704} & [0.663, 0.743]\\
 & Azure & -- & 0.761 & [0.709, 0.804] & 0.161 & [0.129, 0.196]\\
 & aesopocean & S & \textbf{0.872} & [0.842, 0.894] & 0.598 & [0.563, 0.631]\\
 & aesopocean & H & \textbf{0.872} & [0.841, 0.894] & 0.607 & [0.572, 0.640]\\
 & koreanocean & S & 0.853 & [0.816, 0.882] & \textbf{0.821} & [0.789, 0.847]\\
\midrule
Fluency & \textbf{Human consensus} & -- & \textbf{0.807} & [0.768, 0.840] & \textbf{0.739} & [0.695, 0.777]\\
 & Azure & -- & 0.633 & [0.556, 0.698] & 0.078 & [0.056, 0.107]\\
 & aesopkoreanocean & D & \textbf{0.877} & [0.846, 0.902] & 0.577 & [0.539, 0.612]\\
 & koreanaesop & H & \textbf{0.878} & [0.848, 0.901] & 0.469 & [0.431, 0.502]\\
 & koreanocean & H & 0.838 & [0.799, 0.870] & \textbf{0.744} & [0.702, 0.779]\\
\midrule
Prosody & \textbf{Human consensus} & -- & \textbf{0.807} & [0.768, 0.842] & \textbf{0.742} & [0.701, 0.778]\\
 & Azure & -- & 0.737 & [0.672, 0.792] & 0.370 & [0.318, 0.421]\\
 & koreanaesop & S & \textbf{0.883} & [0.852, 0.906] & 0.431 & [0.394, 0.465]\\
 & koreanaesop & D & \textbf{0.883} & [0.852, 0.906] & 0.431 & [0.394, 0.465]\\
 & koreanaesop & H & \textbf{0.883} & [0.853, 0.906] & 0.433 & [0.397, 0.468]\\
 & koreanocean & H & 0.840 & [0.802, 0.872] & \textbf{0.675} & [0.632, 0.712]\\
\bottomrule
\end{tabular}
\par\smallskip
\begin{minipage}{\textwidth}
\footnotesize
\textit{Note.} Statistics average agreement against the three two-human reference pairs; humans are compared with the other two raters. S: construct-specific development Spearman; D: development accuracy MAE, standardized grid; H: development accuracy MAE, normalized historical grid. Each construct includes every configuration within .001 of its highest unrounded Spearman or QWK among 21 combinations; bold marks those metric values and the human baseline. Leaders and near-ties are descriptive, not prospectively selected or statistically established. No displayed leader required clipping. CIs are marginal 95\% percentile-bootstrap intervals (2,000 draws, seed 42), not selection-adjusted. Full results are in the supplement. Paired tests cover all 21 configurations.
\end{minipage}
\end{table*}

\section{Discussion}
Scorer selection depends on construct and intended use: public-development selection prioritizes rank correlation, whereas external QWK favors koreanocean with construct-specific selection rules. These external preferences are descriptive, not independently validated deployment choices. Azure scores average 23.63, 27.88, and 10.57 points above the three-human means for accuracy, fluency, and prosody, respectively; these offsets help explain low QWK despite positive rank associations. No correction was fitted to remove those offsets. XEUS alone approaches the retained combinations on public-development macro correlation (gaps of 0.0017 for accuracy, 0.0016 for fluency, and 0.0041 for prosody in the selected families), so the sweep does not establish a substantial benefit from every additional encoder. A lighter XEUS-only system requires separate external evaluation. Korean-only fluency and prosody share features and targets, so their predictions do not constitute independent learning of two constructs. The human comparison distinguishes the constructs: accuracy QWK exceeds the individual-human baseline with statistical support, fluency QWK is comparable in the observed results, and prosody QWK remains lower. These comparisons measure agreement with the same two-rater references, rather than agreement with an error-free ground truth. Agreement with human ratings supports a particular score interpretation, rather than every educational use of the scores \cite{kane}. The software, trained scoring weights, model configurations, and reproduction materials are available at \url{https://github.com/cburst/oollim}, subject to the licenses documented in the repository. 

\section{Limitations}
The evaluation supports conclusions about one passage and three university EFL settings. Recording-level scores permit paired analysis, but URL-based grouping assumes one recording per speaker. The cohort excludes one long recording, and the 120-second operational limit constrains input coverage. Score-scale mappings do not establish equivalent constructs: the Korean Learners' L2 Speech Corpus (English) supplies shared prosody/fluency supervision, so differences among training families cannot be attributed solely to additional data. Speaker-disjoint fitting also cannot resolve possible upstream exposure of PhoneticXeus to learner corpora. The Azure service revision and remaining request options are unavailable, limiting exact replication of that comparator. The public-development ablation used a recorded selection rule, but neither bootstrap intervals nor contrast adjustment account for the entire model-search process. Broader claims require additional passages, populations, and intended uses \cite{kane}.

\section{Conclusion}
XLS-R, XEUS, and PhoneticXeus with construct-specific ridge scoring support local evaluation of L2 English accuracy, fluency, and prosody on consumer hardware. Every evaluated family exceeds archival Azure correlations across all three constructs. Different leaders emerge for numerical agreement, supporting the reporting of both metrics. 

\section{Generative AI Use Disclosure}
OpenAI Codex assisted with software development, experiment execution, and manuscript drafting and revision; the author is responsible for the submitted content and its verification.

\begin{thebibliography}{99}
\bibitem{deterding} \href{https://doi.org/10.1002/9781118346952.ch4}{D. Deterding. Segmentals. In M. Reed and J. M. Levis, eds., \textit{The Handbook of English Pronunciation}, pp. 67--84. Wiley, 2015.}
\bibitem{cai} \href{https://doi.org/10.1111/lang.70000}{D. Cai, B. Naismith, M. Kostromitina, Z. Teng, K. P. Yancey, and G. T. LaFlair. Developing an Automatic Pronunciation Scorer: Aligning Speech Evaluation Models and Applied Linguistics Constructs. \textit{Language Learning}, 75(S1):170--203, 2025.}
\bibitem{kormos} \href{https://doi.org/10.1016/j.system.2004.01.001}{J. Kormos and M. D\'enes. Exploring measures and perceptions of fluency in the speech of second language learners. \textit{System}, 32(2):145--164, 2004.}
\bibitem{bosker} \href{https://doi.org/10.1177/0265532212455394}{H. R. Bosker, A.-F. Pinget, H. Quen\'e, T. Sanders, and N. H. de Jong. What makes speech sound fluent? The contributions of pauses, speed and repairs. \textit{Language Testing}, 30(2):159--175, 2013.}
\bibitem{cole} \href{https://doi.org/10.1080/23273798.2014.963130}{J. Cole. Prosody in context: A review. \textit{Language, Cognition and Neuroscience}, 30(1--2):1--31, 2015.}
\bibitem{cutler} \href{https://doi.org/10.1002/9781118346952.ch6}{A. Cutler. Lexical Stress in English Pronunciation. In M. Reed and J. M. Levis, eds., \textit{The Handbook of English Pronunciation}, pp. 106--124. Wiley, 2015.}
\bibitem{isaacs} \href{https://doi.org/10.1080/15434303.2013.769545}{T. Isaacs and R. I. Thomson. Rater Experience, Rating Scale Length, and Judgments of L2 Pronunciation: Revisiting Research Conventions. \textit{Language Assessment Quarterly}, 10(2):135--159, 2013.}
\bibitem{munro} \href{https://doi.org/10.1111/j.1467-1770.1995.tb00963.x}{M. J. Munro and T. M. Derwing. Foreign Accent, Comprehensibility, and Intelligibility in the Speech of Second Language Learners. \textit{Language Learning}, 45(1):73--97, 1995.}
\bibitem{segalowitz} N. Segalowitz. \textit{Cognitive Bases of Second Language Fluency}. Routledge, 2010.
\bibitem{hahn} \href{https://doi.org/10.2307/3588378}{L. D. Hahn. Primary Stress and Intelligibility: Research to Motivate the Teaching of Suprasegmentals. \textit{TESOL Quarterly}, 38(2):201--223, 2004.}
\bibitem{azure} Microsoft. Use pronunciation assessment. \textit{Microsoft Learn}. \url{https://learn.microsoft.com/azure/ai-services/speech-service/how-to-pronunciation-assessment}. Accessed October 18, 2026.
\bibitem{williamson} \href{https://doi.org/10.1111/j.1745-3992.2011.00223.x}{D. M. Williamson, X. Xi, and F. J. Breyer. A Framework for Evaluation and Use of Automated Scoring. \textit{Educational Measurement: Issues and Practice}, 31(1):2--13, 2012.}
\bibitem{speechrater} \href{https://doi.org/10.1002/ets2.12198}{L. Chen et al. Automated Scoring of Nonnative Speech Using the SpeechRater v. 5.0 Engine. \textit{ETS Research Report Series}, 2018(1):1--31, 2018.}
\bibitem{gopt} Y. Gong, Z. Chen, I.-H. Chu, P. Chang, and J. Glass. Transformer-Based Multi-Aspect Multi-Granularity Non-Native English Speaker Pronunciation Assessment. \textit{ICASSP}, 2022. \href{https://arxiv.org/abs/2205.03432}{arXiv:2205.03432}.
\bibitem{speechocean} J. Zhang et al. speechocean762: An Open-Source Non-Native English Speech Corpus for Pronunciation Assessment. Interspeech, 2021. \url{https://www.openslr.org/101/}.
\bibitem{kim} \href{https://doi.org/10.21437/Interspeech.2022-10245}{E. Kim, J.-J. Jeon, H. Seo, and H. Kim. Automatic Pronunciation Assessment using Self-Supervised Speech Representation Learning. In \textit{Proc. Interspeech}, pp. 1411--1415, 2022.}
\bibitem{saito} \href{https://doi.org/10.1017/S0272263122000080}{K. Saito, K. Macmillan, M. Kachlicka, T. Kunihara, and N. Minematsu. Automated assessment of second language comprehensibility: Review, training, validation, and generalization studies. \textit{Studies in Second Language Acquisition}, 45(1):234--263, 2023.}
\bibitem{xlsr} \href{https://doi.org/10.21437/Interspeech.2022-143}{A. Babu et al. XLS-R: Self-supervised Cross-lingual Speech Representation Learning at Scale. In \textit{Proc. Interspeech}, pp. 2278--2282, 2022.}
\bibitem{xeus} W. Chen et al. Towards Robust Speech Representation Learning for Thousands of Languages. 2024. \url{https://arxiv.org/abs/2407.00837}.
\bibitem{phonetic} S. Bharadwaj et al. An Empirical Recipe for Universal Phone Recognition. 2026. \url{https://arxiv.org/abs/2603.29042}.
\bibitem{ridge} \href{https://doi.org/10.1080/00401706.1970.10488634}{A. E. Hoerl and R. W. Kennard. Ridge Regression: Biased Estimation for Nonorthogonal Problems. \textit{Technometrics}, 12(1):55--67, 1970.}
\bibitem{ffmpeg} FFmpeg Developers. \textit{FFmpeg Documentation}. \url{https://ffmpeg.org/documentation.html}. Accessed October 18, 2026.
\bibitem{pooling} \href{https://doi.org/10.21437/Interspeech.2017-620}{D. Snyder, D. Garcia-Romero, D. Povey, and S. Khudanpur. Deep Neural Network Embeddings for Text-Independent Speaker Verification. In \textit{Proc. Interspeech}, pp. 999--1003, 2017.}
\bibitem{cmudict} CMU Sphinx. \textit{CMU Pronouncing Dictionary}. \url{https://github.com/cmusphinx/cmudict}. Accessed October 18, 2026.
\bibitem{cmutools} CMU Sphinx. \textit{CMU Pronunciation Dictionary Tools: Phone Tables and Phoneme Sets}. \url{https://github.com/cmusphinx/cmudict-tools}. Accessed October 18, 2026.
\bibitem{ctc} \href{https://doi.org/10.1145/1143844.1143891}{A. Graves, S. Fern\'andez, F. Gomez, and J. Schmidhuber. Connectionist Temporal Classification: Labelling Unsegmented Sequence Data with Recurrent Neural Networks. In \textit{Proc. ICML}, pp. 369--376, 2006.}
\bibitem{shannon} \href{https://doi.org/10.1002/j.1538-7305.1948.tb01338.x}{C. E. Shannon. A Mathematical Theory of Communication. \textit{The Bell System Technical Journal}, 27(3):379--423, 1948.}
\bibitem{editdistance} \href{https://doi.org/10.1145/321796.321811}{R. A. Wagner and M. J. Fischer. The String-to-String Correction Problem. \textit{Journal of the ACM}, 21(1):168--173, 1974.}
\bibitem{han} S. Han, M. Chung, and S. Kim. \href{https://aclanthology.org/2024.lrec-main.334/}{Constructing Korean Learners' L2 Speech Corpus of Seven Languages for Automatic Pronunciation Assessment. In \textit{Proc. LREC-COLING}, pp. 3772--3781, 2024.}
\bibitem{aihub} National Information Society Agency. Korean Learners' L2 Speech Corpus (English). AIHub dataset 71463. \url{https://aihub.or.kr/aihubdata/data/view.do?dataSetSn=71463}.
\bibitem{aesop} K. Yazawa, T. Konishi, and M. Kondo. The J-AESOP corpus: Design, application, and future directions of a Japanese-English bilingual speech corpus. Acoustical Science and Technology, 46(5):553--563, 2025. \url{https://doi.org/10.1250/ast.e25.06}.
\bibitem{cawley} G. C. Cawley and N. L. C. Talbot. On Over-fitting in Model Selection and Subsequent Selection Bias in Performance Evaluation. \textit{Journal of Machine Learning Research}, 11:2079--2107, 2010.
\bibitem{dorans} N. J. Dorans. \href{https://doi.org/10.1002/9781118489772.ch19}{Scores, Scales, and Score Linking. In P. Irwing, T. Booth, and D. J. Hughes, eds., \textit{The Wiley Handbook of Psychometric Testing}. Wiley, 2018.}
\bibitem{esl} \href{https://doi.org/10.1007/978-0-387-84858-7}{T. Hastie, R. Tibshirani, and J. Friedman. \textit{The Elements of Statistical Learning}, 2nd ed. Springer, 2009, secs. 3.4 and 7.}
\bibitem{stemler} S. E. Stemler. \href{https://doi.org/10.7275/zrjx-8p89}{A Comparison of Consensus, Consistency, and Measurement Approaches to Estimating Interrater Reliability. \textit{Practical Assessment, Research \& Evaluation}, 9(4), 2004.}
\bibitem{spearman} \href{https://doi.org/10.2307/1412159}{C. Spearman. The Proof and Measurement of Association between Two Things. \textit{The American Journal of Psychology}, 15(1):72--101, 1904.}
\bibitem{cohen} \href{https://doi.org/10.1037/h0026256}{J. Cohen. Weighted kappa: Nominal scale agreement with provision for scaled disagreement or partial credit. \textit{Psychological Bulletin}, 70(4):213--220, 1968.}
\bibitem{efron} \href{https://doi.org/10.1214/aos/1176344552}{B. Efron. Bootstrap Methods: Another Look at the Jackknife. \textit{The Annals of Statistics}, 7(1):1--26, 1979.}
\bibitem{field} \href{https://doi.org/10.1111/j.1467-9868.2007.00593.x}{C. A. Field and A. H. Welsh. Bootstrapping clustered data. \textit{Journal of the Royal Statistical Society: Series B}, 69(3):369--390, 2007.}
\bibitem{kane} M. T. Kane. \href{https://doi.org/10.1111/jedm.12000}{Validating the Interpretations and Uses of Test Scores. \textit{Journal of Educational Measurement}, 50(1):1--73, 2013.}
\end{thebibliography}
\end{document}
